\section*{第 \ref{ch:b3:complete} 章——完备空间：
贝尔、阿斯科利、斯通--魏尔斯特拉斯}

\begin{solution}{exo:b3:complete:1}
(a) 设 $m \geq n$，则 $f_n - f_m$ 在某个区间之外为零，该区间的
长为 $\frac1n$，且以 $1$ 为界：$\norm{f_n - f_m}_1
\leq \frac1n \to 0$：柯西列。若 $f_n \to f$ 按
$\norm\cdot_1$ 意义下收敛，且 $f$ 连续：对固定的 $\alpha <
\frac12$，有 $\int_0^{\alpha}\abs f = \lim \int_0^\alpha\abs{f -
f_n + f_n} \leq \lim\bigl(\norm{f - f_n}_1 + 0\bigr) = 0$
（$f_n \equiv 0$ 对充分大的 $n$ 成立），故 $f \equiv 0$ 于
$[0, \frac12)$（由连续性）；同理 $f \equiv 1$ 于
$(\frac12, 1]$：没有连续函数能同时满足这两者。因此该空间
不完备——其\omterm{thm:b3:complete:completion}{完备化}为 $L^1$，
在\cref{ch:b3:lp}中构造。

(b) 设 $(x_n)$ 为柯西列；取 $n_1 < n_2 < \cdots$，使得
$\norm{x_{n_{k+1}} - x_{n_k}} \leq 2^{-k}$。级数
$\sum_k (x_{n_{k+1}} - x_{n_k})$ 绝对收敛，从而
收敛；其部分和为 $x_{n_{k+1}} - x_{n_1}$，故
$(x_{n_k})$ 收敛，于是
有收敛子列的柯西列收敛。
\end{solution}

\begin{solution}{exo:b3:complete:2}
(a) 设 $E$ 完备，且具有代数基 $(e_n)_{n\in\N}$，
令 $F_n = \operatorname{Vect}(e_1, \dots, e_n)$，它为闭集
（有限维子空间完备，故闭——
本科二年级），\omterm{def:b3:topology:interior}{内部}为空：若 $B(x, r) \subseteq F_n$，取
$v \notin F_n$，则 $x + \frac{r}{2\norm v}v \in B(x,r)
\setminus F_n$，矛盾。但 $E = \bigcup F_n$（每个向量都是
有限线性组合），与贝尔纲定理矛盾
（见\cref{thm:b3:complete:baire}）。空间 $\R[X]$ 具有
可数基 $(X^n)$，因此没有范数能使其完备。

(b) 固定 $\delta > 0$ 及非空开球 $B$；我们在
$B$ 中找到一点，它属于 $\Omega_\delta = \{x : \operatorname{osc}_xf <
\delta\}$，其中 $\operatorname{osc}_xf = \inf_{V \ni
x}\operatorname{diam} f(V)$。（$\Omega_\delta$ 是\omterm{def:b3:topology:topology}{开集}：
若 $\operatorname{diam}f(V) < \delta$ 对某\omterm{def:b3:topology:topology}{开集} $V \ni x$ 成立，则每
个 $y \in V$ 的振幅均 $< \delta$。）下列集合
\[
F_N = \bigl\{x : \abs{f_n(x) - f_m(x)} \leq \tfrac\delta3\
\ \forall\, n, m \geq N \bigr\}
\]
均为闭集（闭集原像的交），且
覆盖 $\R$（逐点收敛使 $(f_n(x))$ 成为柯西列）。
在完备的 $\bar B$ 中应用贝尔纲定理，存在某个 $F_N \cap
\bar B$ 包含球 $B' = B(x_0, \rho)$。令 $m \to
\infty$ 时，$\abs{f_N - f} \leq \frac\delta3$ 在 $B'$ 上成立。
由 $f_N$ 在 $x_0$ 处的连续性，缩小至 $B'' \ni x_0$，使
$\abs{f_N - f_N(x_0)} \leq \frac\delta3$，则对 $x \in B''$，
\[
\abs{f(x) - f(x_0)} \leq \abs{f - f_N}(x) + \abs{f_N(x) -
f_N(x_0)} + \abs{f_N - f}(x_0) \leq \delta:
\]
$\operatorname{diam} f(B'') \leq 2\delta$，故 $x_0 \in
\Omega_{3\delta} \cap B$。因此每个 $\Omega_\delta$ 皆为\omterm{def:b3:topology:topology}{开集}，且
稠密；$\bigcap_k\Omega_{1/k}$——
即连续点之集——由贝尔纲定理亦稠密。
\end{solution}

\begin{solution}{exo:b3:complete:3}
(a) $f(x) = \frac12(x + \frac ax)$ 把 $[\sqrt a, \infty)$ 映入
自身（算术–几何平均不等式：$f(x) \geq \sqrt{x \cdot \frac ax} = \sqrt a$），
且在该区间上 $f'(x) = \frac12(1 - \frac a{x^2}) \in [0, \frac12)$：
闭（完备）集上的 $\frac12$-利普希茨压缩映射。
不动点：$x = f(x) \iff x^2 = a$，即 $x^* = \sqrt a$。收敛速度
（见\cref{thm:b3:complete:banach}）：$d(x_n, \sqrt a) \leq
2^{-n+1}\,d(x_1, x_0)$。取 $a = 2$，$x_0 = 2$，则 $d(x_1, x_0) =
\frac12$，故 $n = 40$ 即可保证 $2^{-40} < 10^{-12}$。（事实上
牛顿法二次收敛：寥寥数步
迭代即可；压缩估计偏于保守，但
无额外代价。）

(b) $f_m(x) = m + e\sin x$ 是 $e$-利普希茨的，其中 $e < 1$，在
完备的 $\R$ 上有唯一不动点 $x(m)$。对两个参数有：
\[
\abs{x(m) - x(m')} = \abs{f_m(x(m)) - f_{m'}(x(m'))}
\leq \abs{m - m'} + e\abs{x(m) - x(m')},
\]
故 $\abs{x(m) - x(m')} \leq \frac{\abs{m - m'}}{1 - e}$，
从而关于 $m$ 亦为利普希茨的。
\end{solution}

\begin{solution}{exo:b3:complete:4}
(a) 设 $x^*$ 为 $f^p$ 的唯一不动点。则
$f^p(f(x^*)) = f(f^p(x^*)) = f(x^*)$：$f(x^*)$ 为
$f^p$ 的不动点，故 $f(x^*) = x^*$。$f$ 的不动点都是
$f^p$ 的不动点：唯一性由此传递。对 $T$，由归纳法 $\abs{T^pf(x)}
\leq \norm f_\infty x^p/p!$（每次积分各增加一个因子
$\frac xk$），故 $\norm{T^p} \leq a^p/p!$。映射 $S(u) = g +
\lambda Tu$ 满足 $S^p(u) - S^p(v) = \lambda^pT^p(u - v)$，
其范数 $\leq \abs\lambda^pa^p/p!\,\norm{u - v} \to 0$，故某个
$S^p$ 为压缩映射，从而 $S$ 有唯一不动点，即
沃尔泰拉方程 $u = g + \lambda Tu$ 对
\emph{任一} $\lambda$ 皆唯一可解。

(b) $\varphi(x) = d(x, f(x))$ 在紧集 $K$ 上连续：
它在某点 $x_0$ 处取得最小值。若 $f(x_0) \neq x_0$，
则 $\varphi(f(x_0)) = d(f(x_0), f^2(x_0)) < d(x_0, f(x_0)) =
\min\varphi$，矛盾。唯一性：若有两个不动点 $x \ne y$
则 $d(x,y) = d(f(x), f(y)) < d(x,y)$。若无
紧性：$f(x) = x + \frac1x$ 在 $[1, \infty)$ 上满足
$x < y$ 时 $f(y) - f(x) = (y - x)\bigl(1 - \frac1{xy}\bigr) <
y - x$，但 $f(x) > x$ 恒成立，无不动点——严格的距离递减
这一条件弱于一致压缩因子。
\end{solution}

\begin{solution}{exo:b3:complete:5}
(a) 集合 $A = \{g = h\}$ 是对角线
$\Delta_Y$ 在映射 $x \mapsto (g(x), h(x))$ 下的原像，该映射连续；$\Delta_Y$
为闭集，因为 $Y$ 为\omterm{def:b3:topology:hausdorff}{豪斯多夫空间}（对 $y_1 \neq y_2$，
不相交的开\omterm{def:b3:topology:topology}{邻域}可得围绕 $(y_1,
y_2)$ 的开方块，与对角线不交，因此
$\Delta_Y$ 的补集为\omterm{def:b3:topology:topology}{开集}），故 $A$ 为
闭集，包含\omterm{def:b3:topology:interior}{稠密集}，且等于 $X$。用到的是
\cref{thm:b3:complete:extension} 中的唯一性，从而
也即\omterm{thm:b3:complete:completion}{完备化}的唯一性。

(b) $f$ 为等距映射，故一致连续：它可延拓为
$F\colon X \to Y$（见\cref{thm:b3:complete:extension}），仍是
等距的（关系式 $d(F(x), F(x')) = d(x, x')$ 在稠密的点对集上成立，
两边均连续）。同理，
$f^{-1}\colon f(D) \to X$ 可延拓为 $G \colon Y \to X$。复合
映射 $G \circ F$ 连续，且在\omterm{def:b3:topology:interior}{稠密集} $D$ 上恒等，
由 (a) 知它为 $\mathrm{id}_X$；同理 $F \circ G =
\mathrm{id}_Y$。于是 $F$ 为等距双射。\omterm{thm:b3:complete:completion}{完备化}的唯一性：
将结论用于 $D = X$ 在两个
\omterm{thm:b3:complete:completion}{完备化}中稠密的情形。
\end{solution}

\begin{solution}{exo:b3:complete:6}
$\{\sin(nx)\}$：逐点以 $1$ 为界；\emph{不}
\omterm{def:b3:complete:equicontinuous}{等度连续}：在 $x = 0$ 处，$\sin(n\cdot\frac{\pi}{2n}) = 1$
其中 $\frac\pi{2n} \to 0$，故对任何公共的 $\delta$，
在 $\varepsilon = \frac12$ 下均不满足条件。非相对紧（阿斯科利
必要性，\cref{thm:b3:complete:ascoli}）。

$\{x^n\}$：有界；在 $1$ 处不\omterm{def:b3:complete:equicontinuous}{等度连续}：$1 - (1 -
\delta)^n \to 1$（$n \to \infty$ 时，$\delta$ 固定）。非
相对紧——与之相符，其逐点极限
不连续，故任何子列都不一致收敛。

$\{\norm f_\infty \leq 1, \norm{f'}_\infty \leq 1\}$：由中值不等式，
该族为 $1$-利普希茨的，从而\omterm{def:b3:complete:equicontinuous}{等度连续}；有界
故由阿斯科利定理知其相对紧。（不
是紧的：它不是闭的——一致极限未必属于
$\mathcal C^1$ 类；其\omterm{def:b3:topology:interior}{闭包}是 $1$-利普希茨函数中
范数 $\leq 1$ 者。）

$\{\operatorname{Lip}(f) \leq 1, f(0) = 0\}$：\omterm{def:b3:complete:equicontinuous}{等度连续}；
逐点有界（$\abs{f(x)} \leq x \leq 1$）；且在
一致极限下封闭（利普希茨不等式与 $0$ 处的值
在极限下保持）：故紧。
\end{solution}

\begin{solution}{exo:b3:complete:7}
(a) 对 $\norm f_\infty \leq 1$：$\abs{Tf(x)} \leq
\norm k_\infty$，且
\[
\abs{Tf(x) - Tf(x')} \leq \int_0^1\abs{k(x,y) -
k(x',y)}\,\dd y \leq \omega_k\bigl(\abs{x - x'}\bigr),
\]
其中 $\omega_k$ 为 $k$ 在
紧正方形上的一致连续模（海涅）：单位球的像
一致有界且\omterm{def:b3:complete:equicontinuous}{等度连续}，从而相对紧
（阿斯科利）。由线性性，任何有界集的像都相对紧，
故 $T$ 为紧算子。

(b) 子列的陈述即相对紧性定义应用于 $\{Tf_n\}$。若
$T$ 为双射且逆映射连续，则
$T(B(0,1)) \supseteq B(0,
\varepsilon)$ 对某个 $\varepsilon > 0$ 成立（$T^{-1}$ 在
$0$ 处连续）；闭球 $\bar B(0,\varepsilon)$ 为
紧集 $\overline{T(B(0,1))}$ 的闭子集，因而为紧集
——在无穷维的 $\mathcal
C(\intcc01)$ 中不可能，由里斯定理（本科二年级）。
\end{solution}

\begin{solution}{exo:b3:complete:8}
(a) 狄尼：见\cref{lem:b3:complete:dini}——覆盖
论证。(b) 假：移动凸块 $f_n(x) = \max(0, 1 -
\abs{nx - 1})$ 逐点收敛于 $0$（对 $x > 0$ 有 $f_n(x) =
0$，只要 $n > 2/x$；$f_n(0) = 0$），但 $\norm{f_n}_\infty = 1$。
(c) 假：$f_n(x) = x^n$ 递减收敛于不连续函数
$\mathbf 1_{\{1\}}$；$\sup_{[0,1]}\abs{f_n - f} = \sup_{[0,1)}
x^n = 1$。(d) 假：$x^n \downarrow 0$ 在
非紧的 $\intoo01$ 上逐点成立，且 $\sup_{(0,1)}x^n = 1$。每
个狄尼定理的假设都必不可少。
\end{solution}

\begin{solution}{exo:b3:complete:9}
(a) 由线性性，$\int_0^1 fP = 0$ 对每个多项式 $P$ 成立。
取多项式列，使 $P_n \to f$ 一致收敛
（见\cref{cor:b3:complete:weierstrass}）：$\int_0^1 f^2 =
\lim\int_0^1 fP_n = 0$。连续的 $f^2 \geq 0$ 积分为零，
故恒等于零。

(b) 在 $\intcc01$ 上：关于 $x^2$ 的多项式构成含常数的代数，
且分离点（$x \mapsto x^2$ 在 $\intcc01$ 上为单射），
故由斯通–魏尔斯特拉斯定理知其稠密。在 $\intcc{-1}1$ 上：
$x^2$ 在 $\pm x$ 处取相同的值，关于 $x^2$ 的每个
多项式亦然。此类多项式的一致极限为偶
函数。若偶函数 $g_n$ 一致收敛于
恒等映射，则 $x = \lim g_n(x) = \lim g_n(-x) = -x$ 对所有
$x$ 成立：荒谬——故不稠密。分离性假设在
点对 $\{x, -x\}$ 处失效。

(c) 否。对 $f$ 属于代数 $\mathcal A$，该代数由常数与
$\eu^{\iu t}$ 生成，即 $\eu^{\iu nt}$（$n \geq 0$）的线性组合，
则 $\Lambda(f) =
\int_0^{2\pi}f(t)\,\eu^{\iu t}\,\dd t = 0$（每个
$\int_0^{2\pi}\eu^{\iu(n+1)t}\dd t = 0$）。$\Lambda$ 是
关于 $\norm\cdot_\infty$ 连续的（$\abs{\Lambda(f)}\leq
2\pi\norm f_\infty$），从而 $\Lambda$ 在 $\bar{\mathcal
A}$ 上为零；但 $\Lambda(\eu^{-\iu t}) = 2\pi \neq 0$：$\eu^{-\iu t}
\notin \bar{\mathcal A}$。（斯通–魏尔斯特拉斯定理不适用：
$\mathcal A$ 在共轭下不封闭——而
此障碍恰是全纯函数论
将在\cref{ch:b3:holomorphic}中系统化的那一个。）
\end{solution}

\begin{solution}{exo:b3:complete:10}
$F_n = \bigcap_{f \in \mathcal F}\{x : \abs{f(x)} \leq n\}$ 是
闭集之交，故为闭集。逐点有界性
给出 $X = \bigcup_n F_n$。由
\cref{thm:b3:complete:baire}可知，存在 $n_0$ 使得 $U = \mathring F_{n_0} \neq \varnothing$：
在 $U$ 上，$\abs f \leq n_0$ 对每个 $f \in \mathcal F$
同时成立。
\end{solution}

\begin{solution}{exo:b3:complete:11}
(a) 关于 $\norm\cdot_{\mathcal C^1}$ 的柯西列
连同其导数一致柯西：$f_n \to f$
且 $f_n' \to g$ 一致收敛，其中 $f, g$ 连续。
在积分号下取极限（一致收敛允许这样做），由
$f_n(x) = f_n(0) + \int_0^xf_n'(t)\,\dd t$ 可得
$f(x) = f(0) + \int_0^xg$：$f$ 为 $\mathcal C^1$ 函数，且 $f' = g$，
并且 $\norm{f_n - f}_{\mathcal C^1} \to 0$。
完备。

(b) $h(x) = \abs{x - \frac12}$ 是
$\mathcal C^1$ 函数的一致极限，例如 $h_n(x) = \sqrt{(x -
\frac12)^2 + \frac1n}$（$\abs{h_n - h} \leq \frac1{\sqrt n}$
由共轭量估计），但 $h \notin E$：
$\sup$ 范数在 $E$ 上不完备——其\omterm{thm:b3:complete:completion}{完备化}为
$\mathcal C(\intcc01)$。

(c) $f_n(x) = \sin(2\pi nx)$ 满足 $\norm{f_n}_\infty = 1$ 且
$\norm{f_n'}_\infty = 2\pi n \to \infty$：不存在常数 $C$。
（从概念上看：若求导对 sup 范数有界，
则 (a)--(b) 中的两种范数等价，使
$(E, \norm\cdot_\infty)$ 完备——与 (b) 矛盾。
这种无界性正是
一般连续函数可能处处不可导的原因，
见\cref{thm:b3:complete:nowherediff}。）
\end{solution}

\begin{solution}{exo:b3:complete:12}
固定 $\varepsilon > 0$。每个 $F_N$ 是对 $n
\geq N$ 的原像之交，其中闭区间 $\intcc{-\varepsilon}
\varepsilon$ 在\omterm{def:b3:topology:continuity}{连续映射} $x \mapsto f(nx)$ 下的原像，故为闭集。
假设表明每个 $x > 0$ 都落在某个 $F_N$ 中。
对完备的 $\intcc ab$（任取 $0 < a <
b$），应用贝尔纲定理，某个 $F_N$ 在某子区间中稠密；因其为闭集，
故包含区间 $\intcc{a'}{b'}$，其中 $0 < a' < b'$。
于是对每个 $n \geq \max(N, \frac{a'}{b' - a'})$，
区间 $\intcc{na'}{nb'}$ 与 $\intcc{(n+1)a'}{(n+1)b'}$
相互重叠（$nb' \geq (n+1)a'$），故
\[
\bigcup_{n \geq n_0}\intcc{na'}{nb'} \supseteq
\intco{n_0a'}{+\infty},
\]
而每个 $t \geq n_0a'$ 都可写成 $t = nx$，其中 $n \geq N$，$x \in
\intcc{a'}{b'} \subseteq F_N$，故 $\abs{f(t)} \leq \varepsilon$。
因此 $\limsup_{t\to\infty}\abs f \leq \varepsilon$ 对一切
$\varepsilon$ 成立，即 $f \to 0$。需要完整假设，
因为 $F_N$ 必须\emph{覆盖}一段区间上的 $x$
——若只取有理数 $x$，则 $F_N$ 的并为
可数集，贝尔纲定理无能为力；事实上存在
连续反例，沿所有有理射线趋于零
但在无穷远处不趋于零。
\end{solution}

\begin{solution}{pb:b3:complete:1}
\textbf{1.} 对 $k$ 作归纳：若 $\norm{\varphi_n(t_k) - x_0}
\leq M(t_k - t_0) \leq MT \leq b$，则点 $(t_k,
\varphi_n(t_k))$ 属于 $R$，故斜率 $f(t_k,
\varphi_n(t_k))$ 有定义，范数 $\leq M$；于是对 $t \in
[t_k, t_{k+1}]$，有 $\norm{\varphi_n(t) - x_0} \leq
\norm{\varphi_n(t_k) - x_0} + M(t - t_k) \leq M(t - t_0) \leq
MT \leq b$。

\textbf{2.} 每个仿射段的斜率范数 $\leq M$；
斜率以 $M$ 为界的分段仿射函数
是 $M$-利普希茨的（沿分点逐段串接）。

\textbf{3.} 族 $(\varphi_n)$ 逐点有界
（取值于 $\bar B(x_0, b)$ 中），且\omterm{def:b3:complete:equicontinuous}{等度连续}（公共
利普希茨常数为 $M$）。由阿斯科利定理
（见\cref{thm:b3:complete:ascoli}），可取出 $\varphi_{n_j} \to
\varphi$ 在 $[t_0, t_0 + T]$ 上一致收敛。界在极限下保持，
即 $\varphi$ 为 $M$-利普希茨的，$\varphi(t_0) = x_0$。

\textbf{4.} $R$ 紧，$f$ 连续，故一致连续
（海涅，见\cref{cor:b3:topology:heineborel}）；记
$\delta(\varepsilon)$ 为一致连续模。对非分点 $t \in (t_k,
t_{k+1})$ 上：$\varphi_n'(t) = f(t_k, \varphi_n(t_k))$，而
$f$ 的两个求值点在时间上相差 $\abs{t - t_k} \leq T/n$，
在空间上相差 $\norm{\varphi_n(t) - \varphi_n(t_k)} \leq MT/n$。
对 $n \geq n_0$ 且 $(1 + M)T/n_0 < \delta(\varepsilon)$，
有 $\norm{\Delta_n(t)} \leq \varepsilon$。

\textbf{5.} 在每个 $[t_k, t_{k+1}]$ 上 $\varphi_n$ 为仿射函数，故
$\varphi_n(t_{k+1}) - \varphi_n(t_k) = \int_{t_k}^{t_{k+1}}
\varphi_n'(s)\dd s$，其导数即为常值斜率；
对各段求和（末段在 $t$ 处截断）：
$\varphi_n(t) = x_0 + \int_{t_0}^t\varphi_n'(s)\,\dd s$.
写 $\varphi_n' = f(s, \varphi_n(s)) + \Delta_n(s)$，
（被积函数分段连续，跳跃点有限），即得
上式。

\textbf{6.} 给定 $\varepsilon$：当 $j$ 充分大时，
$\norm{\varphi_{n_j} - \varphi}_\infty < \delta(\varepsilon)$，
从而 $\norm{f(s, \varphi_{n_j}(s)) - f(s, \varphi(s))} \leq
\varepsilon$ 对所有 $s$ 成立：被积函数
一致收敛，且 $\int_{t_0}^t f(s,\varphi_{n_j}(s))\dd s \to
\int_{t_0}^tf(s, \varphi(s))\dd s$ 关于 $t$ 一致成立。另有
$\norm{\int_{t_0}^t\Delta_{n_j}} \leq T\sup\norm{\Delta_{n_j}}
\to 0$（第 4 问）。在第 5 问的等式中取极限：
$\varphi(t) = x_0 + \int_{t_0}^tf(s,
\varphi(s))\,\dd s$。

\textbf{7.} 被积函数 $s \mapsto f(s, \varphi(s))$ 连续，
故右端是 $\mathcal C^1$ 函数（关于 $t$），
其导数为 $f(t, \varphi(t))$，即 $\varphi$ 在
$[t_0, t_0 + T]$ 上解柯西问题。对左半部分，令 $g(t, x) =
-f(2t_0 - t, x)$，它在反射后的矩形上连续，并具有
相同界 $M$；若 $\psi$ 为 $y' = g(t,y)$ 的解，
且 $y(t_0) = x_0$（于 $[t_0, t_0 + T]$ 上），则 $\varphi(t) =
\psi(2t_0 - t)$ 为原方程在 $[t_0 - T,
t_0]$ 上的解；两半拼接成一个 $\mathcal C^1$ 解（两侧
在 $t_0$ 处的单侧导数均等于 $f(t_0, x_0)$）。——
\emph{皮亚诺定理}：若 $f$ 连续，则
对每个初值都存在局部解。

\textbf{8.} 连续性是显然的。在 $0$ 附近不满足利普希茨条件：
$\abs{f(h) - f(0)} = 2\sqrt h$，而 $2\sqrt h \leq L h$
在 $h < 4/L^2$ 时不成立。

\textbf{9.} 当 $t \leq c$ 时，$x_c \equiv 0$ 为解。当 $t \geq
c$ 时：$x_c'(t) = 2(t - c) = 2\sqrt{(t-c)^2} = 2\sqrt{\abs{x_c}}$。
在 $t = c$ 处两个单侧导数均为 $0$，故 $x_c$
是 $\mathcal C^1$ 类，整体上为解，且 $x_c(0) = 0$ 对每个
$c \geq 0$ 成立——连同 $x \equiv 0$，
过原点的解组成一个连续统。

\textbf{10.} 皮卡迭代构造 $\Phi(x)(t) = x_0 +
\int_0^t f(x(s))\dd s$，并需要 $\norm{\Phi(x) - \Phi(y)}
\leq k\norm{x - y}$（$k < 1$）在某个合适的球上成立——
这一结论由 $f$ 的利普希茨界导出，并经
积分号传递。$\sqrt\cdot$ 在 $0$ 附近不存在利普希茨界，
无论区间或球如何选取都无法修补它。事实上，
任何关于唯一性的证明都不可能成功：唯一性\emph{不成立}
（第 9 题）。

\textbf{11.} 在 $c_0$ 中，单位向量 $e_n = (0, \dots, 0, 1,
0, \dots)$ 满足 $\norm{e_n - e_m}_\infty = 1$，其中 $n \neq
m$：任何子列都不是柯西列，因而闭单位球不是
紧的。失效的一步是提取子列（第 3 题）：
阿斯科利定理用于 $\mathcal C([t_0, t_0+T], E)$ 时，要求取值
位于一个有界集都相对紧的空间中——
在 $\R^d$ 中成立（波尔查诺--魏尔斯特拉斯），在 $c_0$ 中不成立；逐点
提取子列也已不可用（事实上，迪厄多内的例子
根本没有局部解）。

\textbf{12.} \emph{皮亚诺}：若 $f$ 在
$(t_0, x_0)$ 在 $\R\times\R^d$ 中的\omterm{def:b3:topology:topology}{邻域}上连续 $\Rightarrow$ 存在一个
$\mathcal C^1$ 类解 $x' = f(t,x)$，$x(t_0) = x_0$，定义于
某个 $[t_0 - T, t_0 + T]$ 上。\emph{柯西--利普希茨}
（参见\cref{ch:b3:ode}）：若进一步 $f$ 关于
$x$ 局部利普希茨，则解是\emph{唯一}的（任意两解在
公共区间上一致）——即存在性与唯一性。下面这一对
$(x' = 2\sqrt{\abs x},\ x(0) = 0)$ 把两个定理区分开来。

\textbf{13.} 当 $\omega(r) = Lr$ 时：$\int_0^1\frac{\dd r}{Lr} =
+\infty$：满足条件。$\omega(r) = r\log\frac1r$（在 $0$ 附近）：
$\int\frac{\dd r}{r\log(1/r)} = \bigl[-\log\log\frac1r\bigr]
\to +\infty$（$r \to 0$ 时）：满足条件——但 $\omega(r)/r =
\log\frac1r \to \infty$：不满足利普希茨条件。$\omega(r) = 2\sqrt
r$：$\int_0^1\frac{\dd r}{2\sqrt r} = \bigl[\sqrt r\bigr]_0^1
= 1 < \infty$：不满足条件。

\textbf{14.} 将两个积分形式相减：
\[
x_1(t) - x_2(t) = x_1(s) - x_2(s) + \int_s^t\bigl(f(v,
x_1(v)) - f(v, x_2(v))\bigr)\dd v,
\]
取范数，并利用奥斯古德模：$\delta(t) \leq
\delta(s) + \int_s^t\omega(\delta(v))\,\dd v$。

\textbf{15.} $\tau$ 有定义（$\delta(t_0) = 0$），且
$\delta(\tau) = 0$（由连续性），且 $\delta > 0$ 在
$\intoc\tau{t_1}$ 上成立。取定 $s \in \intoo\tau{t_1}$，则 $u$
是 $\mathcal C^1$ 类，$u(s) = \delta(s) > 0$，且 $u \geq \delta$ 于
$[s, t_1]$ 上（第 14 题），且 $u' = \omega(\delta) \leq
\omega(u)$（$\omega$ 单调不减，$u > 0$）。两边相除并
积分：
\[
\int_{\delta(s)}^{u(t_1)}\frac{\dd r}{\omega(r)}
= \int_s^{t_1}\frac{u'(v)}{\omega(u(v))}\,\dd v \leq t_1 - s
\leq t_1 - \tau .
\]
尽管 $u$ 依赖于 $s$，但 $u(t_1) \geq
\delta(t_1)$ 对每个 $s$ 都成立，故左端至少为
$\int_{\delta(s)}^{\delta(t_1)}\frac{\dd r}{\omega(r)}$，
它趋于 $+\infty$（当 $s \downarrow \tau$ 时，
此时 $\delta(s) \to \delta(\tau) = 0$，且积分在
$0$ 处发散）：与右端有界矛盾。故
$\delta \equiv 0$：唯一性成立。

\textbf{16.} 利普希茨情形即 $\omega = Lr$，唯一性得以恢复。
对 $x' = x\log\frac1{\abs x}$，右端满足
奥斯古德模 $\omega(r) = r\log\frac1r$（在 $0$ 附近；由中值不等式
作用于 $x \mapsto x\log\frac1x$ 上，其导数
$\log\frac1x - 1$ 无界——利普希茨条件不成立，奥斯古德条件成立）：
解唯一；注意 $x \equiv 0$ 是一个解，故其他解不能
触及 $0$。对 $\omega(r) = 2\sqrt r$，
$\int_0^{h^2}\frac{\dd r}{2\sqrt r} = h$——“奥斯古德
预算”自 $0$ 攀升到高度 $h^2$ 所需的时间恰为
$h$，事实上 $x_c(c + h) = h^2$：解用去的时间
$h$ 恰好用于收敛积分所允许的事。
积分发散，即是无法在
有限时间内离开 $0$；收敛则是逃逸的途径。

\textbf{17.} $G(t) = \int_{t_0}^t(Le(s) + \eta(s))\dd s$ 是
$\mathcal C^1$ 类，且 $G' = Le + \eta \leq LG + \sup\eta$
（由假设 $e \leq G$）。于是 $\bigl(\eu^{-Lt}(G +
\tfrac{\sup\eta}L)\bigr)' = \eu^{-Lt}(G' - LG - \sup\eta)
\leq 0$：括号项递减，故 $G(t) + \frac{\sup\eta}L
\leq \eu^{L(t - t_0)}\bigl(G(t_0) + \frac{\sup\eta}L\bigr) =
\eu^{L(t-t_0)}\frac{\sup\eta}L$，即 $e(t) \leq G(t) \leq
\frac{\sup\eta}L(\eu^{L(t-t_0)} - 1)$。

\textbf{18.} 定量偏差：在 $(t_k, t_{k+1})$ 上，
$\Delta_n(t) = f(t_k, \varphi_n(t_k)) - f(t, \varphi_n(t))$
其中 $\abs{t - t_k} \leq T/n$，且 $\norm{\varphi_n(t) -
\varphi_n(t_k)} \leq MT/n$，故 $\norm{\Delta_n} \leq L'T/n +
LMT/n$。将 $\varphi_n$ 的积分恒等式
（第 5 题）与 $\varphi$ 的作差，并利用利普希茨界：
\[
\norm{\varphi_n(t) - \varphi(t)} \leq
\int_{t_0}^t L\,\norm{\varphi_n - \varphi}(s)\,\dd s +
\int_{t_0}^t\norm{\Delta_n(s)}\,\dd s,
\]
再由第 17 题（取 $\eta = \norm{\Delta_n}$）即得
所述的 $O(1/n)$ 界，$\varphi$ 的唯一性由
柯西--利普希茨（或奥斯古德）定理保证。即便不谈速率，
等度有界、等度利普希茨的 $(\varphi_n)$ 的每个子列都有子列
收敛（由阿斯科利定理及第二部分）到\emph{某个}解，
而由唯一性，该解必为 $\varphi$：
一个数列，若其每个子列都有子列
收敛于同一极限，则该数列收敛。

\textbf{19.} 由 $\varphi_n(t_k) = 0$，斜率 $2\sqrt0 =
0$ 得 $\varphi_n(t_{k+1}) = 0$；由归纳法 $\varphi_n
\equiv 0$，收敛于零解。对
$\varepsilon > 0$，在 $[\varepsilon', \infty)$ 上，取
$\varepsilon' < \varepsilon$ 时，函数 $2\sqrt x$ 满足
利普希茨条件，故第 18 题适用，欧拉格式收敛到
过 $(0, \varepsilon)$ 的唯一解，即 $x(t) =
(t + \sqrt\varepsilon)^2$（验证：$x' = 2(t + \sqrt
\varepsilon) = 2\sqrt x$）。当 $\varepsilon \to 0$ 时，$(t +
\sqrt\varepsilon)^2 \to t^2$ 在 $[0,1]$ 上一致成立，于是二重
极限为 $x_0(t) = t^2$，而非 $0$。任意小的
初值扰动即可使格式
由一个解转向另一个解：将非唯一性
理解为数值不稳定。

\textbf{20.} 非空性：由第二部分。每个解满足
$\norm{x'} = \norm{f(t, x)} \leq M$：$\mathcal S$ 是
一致 $M$ 利普希茨且一致有界的（取值于
$\bar B(x_0, b)$ 中）。闭性：若 $x_n \in \mathcal S \to x$
一致，在 $x_n(t) = x_0 +
\int_{t_0}^tf(s, x_n(s))\dd s$ 中取极限（被积函数由
$f$ 在紧集 $R$ 上的一致连续性而一致收敛）：
$x \in \mathcal S$。阿斯科利：$\mathcal S$ 是闭的，
有界且\omterm{def:b3:complete:equicontinuous}{等度连续}的 $\mathcal C([t_0, t_0+T],
\R^d)$ 的子集：紧。

\textbf{21.} 任一解单调不减（$x' = 2\sqrt{\abs
x} \geq 0$），且 $x(0) = 0$，故 $x \geq 0$。设 $c =
\sup\{t \in [0,1] : x(t) = 0\}$（可能 $c = \infty$，若
$x \equiv 0$，此时 $x = x_\infty$）。对 $t > c$，$x
> 0$（由单调性及 $c$ 的定义），且在此处 $(\sqrt
x)' = \frac{x'}{2\sqrt x} = 1$，故 $\sqrt{x(t)} = t - c$
（由 $c$ 处的连续性）：$x = x_c$。映射 $c \mapsto
x_c$ 的连续性：对 $c, c' \in [0, 1]$，$\sup_{[0,1]}\abs{(t - c)_+^2 -
(t - c')_+^2} \leq 2\abs{c - c'}$（映射 $c \mapsto
(t-c)_+^2$ 是 $2$-利普希茨的（关于 $t \in [0,1]$ 一致），且
由于 $x_c \equiv 0$ 于 $[0,1]$ 上（对每个 $c \geq 1$），故
族约化为 $\mathcal S = \{x_c : c \in [0,1]\}$（其中
$x_1 = 0 = x_\infty$）。故 $\mathcal S$ 是
紧连通的区间 $[0,1]$ 在\omterm{def:b3:topology:continuity}{连续映射} $c
\mapsto x_c$ 下的像，紧且连通。解的
“漏斗”是一段连续线段，自 $t^2$（立即
逃逸）降至 $0$（永久静止）。

\textbf{22.} 求值映射 $\operatorname{ev}_t\colon
\mathcal C([0,1]) \to \R$，$x \mapsto x(t)$，是连续的
（$\abs{x(t) - y(t)} \leq \norm{x - y}_\infty$），故 $\mathcal
S(t) = \operatorname{ev}_t(\mathcal S)$ 为紧集的连续像，
故紧——又为连通集的连续像，故连通。
对于例子：$x_c(t) = (t - c)_+^2$ 随 $c$ 变动，
取遍 $[0, 1]$ 时，连续地扫过从 $t^2$（$c = 0$ 时）降至
$0$（$c \geq t$ 时），连续：$\mathcal S(t) =
\intcc0{t^2}$。在每一时刻，漏斗的截面是
一条完整线段：在静止与最大逃逸之间，每种
折中都由某个真实解实现。

\textbf{23.} 将积分形式 $x(t) = x_0 +
\int_{t_0}^tf(s, x(s))\dd s$ 及 $y$ 的类似式相减，并
令 $e(t) = \norm{x(t) - y(t)}$，$d = \norm{x_0 - y_0}$：
\[
e(t) \leq d + \int_{t_0}^tL\,e(s)\dd s = G(t).
\]
于是 $G(t_0) = d$，且 $G' = Le \leq LG$，故 $\bigl(\eu^{-L(t
- t_0)}G\bigr)' \leq 0$，且 $G(t) \leq d\,\eu^{L(t - t_0)}$；
从而 $e(t) \leq d\,\eu^{L(t-t_0)}$。关于最优性：对 $x' =
Lx$，过 $x_0$ 与 $y_0$ 的解分别为 $x_0\eu^{L(t -
t_0)}$ 与 $y_0\eu^{L(t-t_0)}$，二者距离恰为
$d\,\eu^{L(t-t_0)}$。当 $d = 0$ 时，$e \equiv 0$：
柯西--利普希茨唯一性，两行即可重新推出。又
对固定的 $t \in [t_0, t_0 + T]$，$\norm{x(t; x_0) - x(t;
y_0)} \leq \eu^{LT}\norm{x_0 - y_0}$：流关于初始条件
是利普希茨的——确定性依赖，
代价为可控的指数因子。

\textbf{24.} 在 $\intoc01$ 上，$\Omega$ 有定义（
积分在 $0$ 处由假设收敛），$\mathcal C^1$ 且
$\Omega' = \frac1\omega > 0$，故为递增双射，映到
$\intoc0{\Omega(1)}$ 上，且 $\Omega(x) \to 0$（$x \to 0^+$）。
其逆 $g \colon \intoc0{\Omega(1)} \to \intoc01$ 是
$\mathcal C^1$ 的，并满足
\[
g'(t) = \frac1{\Omega'(g(t))} = \omega\bigl(g(t)\bigr) > 0,
\qquad g(t) \xrightarrow[t \to 0^+]{} 0 .
\]
将 $g$ 以 $0$ 延拓至 $t \leq 0$：连续性显然，且在
$t = 0$ 处，对 $t > 0$，
\[
\frac{g(t)}t = \frac1t\int_0^tg'(s)\dd s
= \frac1t\int_0^t\omega\bigl(g(s)\bigr)\dd s
\leq \omega\bigl(g(t)\bigr) \xrightarrow[t\to0^+]{} 0
\]
（$\omega$ 单调不减，$g$ 递增，$\omega(0^+) = 0$）：
$g'(0) = 0 = f(g(0))$，且 $g' = f(g)$ 在
$0$ 两侧均成立。故 $x \equiv 0$ 与 $g$ 是两个不同的 $\mathcal
C^1$ 解，过 $(0, 0)$：当 $\int_0\frac{\dd r}
{\omega(r)}$ 收敛时，唯一性失效——第 15 题的
发散恰为其边界。对 $\omega(r) = 2\sqrt
r$：$\Omega(x) = \int_0^x\frac{\dd r}{2\sqrt r} = \sqrt x$，
$g(t) = t^2$，时间平移即得整族
第三部分中的 $x_c$。

\textbf{25.} 取步长 $h = \frac1n$：$\varphi_n(t_{k+1}) =
\varphi_n(t_k) + h\,\varphi_n(t_k) = (1 + h)\varphi_n(t_k)$,
故 $\varphi_n(1) = (1 + \frac1n)^n$，经 $n$ 步。
展开：
\[
n\log\Bigl(1 + \frac1n\Bigr)
= n\Bigl(\frac1n - \frac1{2n^2} + O\Bigl(\frac1{n^3}\Bigr)
\Bigr) = 1 - \frac1{2n} + O\Bigl(\frac1{n^2}\Bigr),
\]
取指数得 $(1 + \frac1n)^n = \eu\,\eu^{-1/(2n) +
O(n^{-2})} = \eu\bigl(1 - \frac1{2n} + O(n^{-2})\bigr)$。于是
$t = 1$ 处的误差为 $\eu - \varphi_n(1) =
\frac{\eu}{2n} + O(n^{-2})$，即第 18 题中的 $O(\frac1n)$，
此处其精确常数为 $\frac\eu2$。数值上，$n =
10$：$1.1^{10} = 2.5937424601$（$1.1^2 = 1.21$，$1.1^4 =
1.4641$，$1.1^8 = 2.14358881$，再乘 $1.21$），且 $\eu -
2.59374 \approx 0.12454$，与渐近预测
$\frac{\eu}{20} \approx 0.13591$ 相比：在
$O(n^{-2})$ 修正量内吻合，其首项在此将
预测值拉向观测值。
\end{solution}
